\documentclass[12pt,a4paper]{article}
\IfFileExists{pt-common-fr.sty}{\usepackage{pt-common-fr}}{\usepackage{../shared/pt-common-fr}}
\usepackage[margin=2.5cm]{geometry}
\usepackage{setspace}
\onehalfspacing

% Environnement Prédiction (utilisé en §5 ; partage le compteur de theoreme)
\newtheorem{prediction}[theoreme]{Prédiction}

\title{\textbf{Holonomie sans potentiel :\\
la Théorie de la Persistance et le débat\\
Aharonov--Bohm / non-localité}}
\author{Yan Senez\\
\texttt{yan.senez@protonmail.com}}
\date{Mai 2026}

\begin{document}
\maketitle

\begin{abstract}
Le débat fondationnel opposant la réalité du potentiel vecteur
(position de Bohm--Hiley), la non-localité intrinsèque du champ
(lectures fortes d'EPR/Bell), et la primauté de l'holonomie
(position de Wu--Yang--Healey) n'admet, dans le cadre classique de
la théorie des champs sur variété, aucune prédiction expérimentale
qui le tranche. La Théorie de la Persistance (PT), dérivée
intégralement à partir du paramètre de symétrie $s = 1/2$, fait
émerger cette résolution à partir d'un théorème structurel et
d'une prédiction falsifiable.

Nous démontrons que (i) le champ dynamique fondamental est unique
(la suite des écarts entre premiers consécutifs, BA0/T0) ;
(ii) le couple potentiel-champ classique $A_\mu \leftrightarrow
F_{\mu\nu}$ est entièrement dérivé du gradient et du Hessien du
potentiel scalaire de persistance $S = -\ln \alpha$, donc
secondaire ; (iii) l'objet primitif est l'angle d'holonomie
modulaire $\theta_p$, jauge-invariant et topologique ; et (iv) la
corrélation $I(\Phi_p\,;\,\Phi_q\,|\,G(\mu))$ entre canaux
modulaires distincts est invariante sous la famille géométrique
émergente $G(\mu)$ (théorème démontré dans l'article companion
\cite{PT_CRT_DECOUPLING}). Cette invariance reclassifie la
non-localité spatiale apparente de Bell--EPR comme une corrélation
arithmétique CRT intra-canal, et la phase Aharonov--Bohm comme un
invariant topologique du canal $\theta_3$ indépendant de la
métrique. La prédiction QH1a (information mutuelle invariante par
géométrie spatiale) est testable sur plateformes d'atomes froids,
ions piégés, et molécules diatomiques.

Les Théorèmes~4.1 et~4.2 ainsi que leur cadre structurel sont
formellement vérifiés en Lean~4 + Mathlib sous le namespace
\texttt{PT.CrtDecoupling} de l'article companion
\cite{PT_CRT_DECOUPLING} (Phases~1 à~4.1, 0~\texttt{sorry},
0~\texttt{axiom}).
\end{abstract}

\tableofcontents

\section{Le débat fondationnel : trois positions classiques}
\label{sec:debat}

Depuis l'article fondateur d'Aharonov et Bohm
\cite{AharonovBohm1959}, l'électrodynamique quantique pose une
question fondationnelle qui résiste à toute réduction à
l'expérience : que doit-on considérer comme l'objet physiquement
réel parmi le 4-potentiel $A_\mu$, le tenseur de champ
$F_{\mu\nu}$, et la phase d'holonomie
$\oint_\gamma A \cdot d\ell$ ? Trois positions se partagent la
littérature contemporaine, sans qu'aucune expérience classique
n'ait pu les départager. Cette section les rappelle dans leur
forme la plus tranchée, dans le seul but de fixer ce que la
Théorie de la Persistance (PT) prétend trancher. La revue n'a pas
vocation à être exhaustive --- pour les développements
philosophiques on renvoie aux monographies de Healey
\cite{Healey2007} et Maudlin~\cite{Maudlin2011}.

\subsection{Position A : la réalité du potentiel vecteur}
\label{ssec:position_A}

La position A, héritée du programme de Bohm--Hiley
\cite{BohmHiley1993}, soutient que le potentiel $A_\mu$ est
physiquement réel et que le tenseur $F_{\mu\nu} = \partial_\mu
A_\nu - \partial_\nu A_\mu$ en est une grandeur dérivée. L'argument
empirique est l'effet Aharonov--Bohm \cite{AharonovBohm1959} :
le déphasage interférométrique
$\Delta\varphi = (e/\hbar) \oint A \cdot d\ell$ enregistré par une
particule qui contourne un solénoïde idéal est mesurable même là
où $E = B = 0$ sur tout son trajet. Aucun champ ne traverse le
trajet ; aucun objet local ne peut donc rendre compte de l'effet.
Le potentiel, lui, est non nul (à jauge près) le long du trajet,
et le programme bohmien y voit la preuve qu'il faut le considérer
comme une grandeur d'ontologie supérieure.

La difficulté connue de la position A est la dépendance de jauge.
Le potentiel $A_\mu$ n'est défini qu'à une transformation
$A_\mu \mapsto A_\mu + \partial_\mu \Lambda$ près, et toute
quantité physique mesurable doit être invariante sous cette
transformation. Or la phase Aharonov--Bohm n'est invariante de
jauge \emph{que dans son intégrale fermée} ; ses valeurs locales
ne le sont pas. La position A doit donc soit accepter une réalité
physique qui n'est pas elle-même invariante de jauge --- ce que
beaucoup considèrent comme une concession ontologique majeure ---
soit traiter $A_\mu$ comme un objet hybride dont seule la classe
d'équivalence sous transformations de jauge serait réelle.

\subsection{Position B : la non-localité intrinsèque du champ}
\label{ssec:position_B}

La position B accepte $F_{\mu\nu}$ comme l'objet réel mais en
fait un objet intrinsèquement non-local. Dans la lecture forte
proposée par Mead \cite{Mead2000} et dans le sillage des analyses
de Bell~\cite{Bell1964} et des tests d'Aspect \cite{Aspect1982} et
Hensen et al.~\cite{Hensen2015}, le champ
électromagnétique n'est pas une distribution sur l'espace-temps
classique : il porte une structure de corrélation entre points
spatialement séparés qui ne s'explique pas par la variation
locale de ses composantes. La phase Aharonov--Bohm devient alors
le signal d'une telle non-localité, et les corrélations EPR celui
de son extension à la matière.

L'objection structurale à la position B est qu'elle exige de
modifier l'ontologie classique du champ pour la rendre
intrinsèquement non-locale, sans préciser quel mécanisme
mathématique implémente cette modification. Les approches
formelles disponibles (fibrés de connexions, théories de jauge,
opérateurs de Wilson) sont toutes interprétables en termes
strictement locaux moyennant un choix adapté de variables. La
non-localité de la position B reste donc, dans la majorité des
formulations, une propriété de l'\emph{interprétation} du champ
plutôt que de sa structure mathématique.

\subsection{Position C : la primauté de l'holonomie}
\label{ssec:position_C}

La position C, due à Wu et Yang \cite{WuYang1975} et reprise par
Healey \cite{Healey2007} et Belot \cite{Belot1998}, soutient
qu'aucun des deux objets locaux $A_\mu, F_{\mu\nu}$ n'est
fondamental : $A_\mu$ parce qu'il dépend de la jauge,
$F_{\mu\nu}$ parce qu'il ne contient pas l'information de la
phase Aharonov--Bohm dans les régions où il s'annule. Le seul
objet à la fois jauge-invariant et sensible à l'expérience
Aharonov--Bohm est la boucle de Wilson
$W_\gamma = \exp\bigl[i (e/\hbar) \oint_\gamma A \cdot d\ell\bigr]$,
fonctionnelle non-locale des chemins fermés~$\gamma$. C'est elle,
et non $A_\mu$ ou $F_{\mu\nu}$, qui mérite le statut d'objet
primitif.

La position C est conceptuellement satisfaisante --- elle
dissout l'asymétrie entre la jauge-dépendance de $A_\mu$ et
l'incomplétude de $F_{\mu\nu}$ --- mais elle laisse ouverte la
question de l'\emph{origine} des boucles de Wilson : pourquoi
ces objets non-locaux particuliers, et non d'autres? La théorie
de jauge classique les introduit comme une construction
mathématique post hoc, à partir d'une connexion qui reste
elle-même primaire dans la formulation. La position C reste donc,
en un sens, une description et non une explication.

\subsection{Le constat opérationnel}
\label{ssec:constat}

Les trois positions A, B, C font les mêmes prédictions
expérimentales dans le cadre de l'électrodynamique quantique
standard. La phase Aharonov--Bohm est mesurable dans les trois
cadres ; les corrélations EPR/Bell sont reproduites dans les
trois ; les sections efficaces et les durées de vie sont
identiques. Aucune mesure n'a pu, à ce jour, départager les trois
positions au niveau de la prédiction quantitative. Le débat est
donc, à l'heure actuelle, un débat d'interprétation.

La Théorie de la Persistance modifie cet état de fait sur deux
points. D'une part, elle promeut la position~C en \emph{théorème}
plutôt qu'en interprétation : l'angle d'holonomie modulaire
$\theta_p$ est l'unique invariant primitif satisfaisant quatre
conditions structurales que ni $A_\mu$ ni $F_{\mu\nu}$ ne peuvent
remplir simultanément (Section~\ref{sec:holonomie_primitive}).
D'autre part, elle produit une prédiction expérimentale --- la
prédiction QH1a de la Section~\ref{sec:prediction_qh1a} --- dont
le verdict empirique distingue effectivement la PT (et avec elle
la position~C) d'une théorie standard sur variété courbe. Le
débat fondationnel n'est plus seulement un débat
d'interprétation : il est, sous PT, un débat tranchable.

\section{Le champ unique et la connexion modulaire de PT}
\label{sec:champ_connexion}

Avant d'évaluer la position C dans le cadre PT
(Section~\ref{sec:holonomie_primitive}), il faut établir trois
faits structurels : que la PT admet un unique degré de liberté
dynamique (\S~\ref{ssec:BA0}), que ce degré de liberté porte
naturellement une connexion de jauge modulaire
(\S~\ref{ssec:connexion_jauge}), et que cette connexion induit
les angles d'holonomie $\theta_p$ par une identité algébrique
explicite (\S~\ref{ssec:T6_holonomie}). Ces trois faits sont
établis dans la monographie~\cite{PT_Monograph} et reproduits ici
sans démonstration.

\subsection{Le champ dynamique unique BA0/T0}
\label{ssec:BA0}

Le théorème de fermeture T0 du chapitre~25 de la
monographie~\cite{PT_Monograph} établit que la suite des écarts
entre nombres premiers consécutifs
\begin{equation}
  g_n \;:=\; p_{n+1} - p_n
  \qquad (n \geq 1)
  \label{eq:g_n_def_fr}
\end{equation}
est l'\emph{unique} degré de liberté dynamique du crible
d'Eratosthène, sous les conditions structurelles U1--U4
(anti-diagonalité de $T_3$, identité GFT, géométrie canonique de
$\mathcal{D}_{\rm KL}$ et Fisher, point fixe $\mustar = 15$). Le
postulat BA0 de la monographie --- selon lequel
$\{g_n\}_{n \geq 1}$ est \emph{le} champ dynamique de la théorie ---
est par conséquent promu de postulat en théorème, et la théorie
admet \emph{zéro} axiome dynamique au-delà des outils standards
de la théorie des nombres.

C'est ce constat qui distingue radicalement la PT des théories
classiques de jauge : alors que ces dernières doivent postuler un
champ vectoriel $A_\mu$ comme structure primitive de
l'électrodynamique, la PT déduit \emph{toute} la structure de
jauge --- y compris l'existence même d'une connexion --- du seul
champ scalaire $\{g_n\}$ et des contraintes de cohérence
arithmétique. La connexion de jauge sera, dans ce qui suit, un
objet \emph{dérivé}, non un postulat.

\subsection{La connexion de jauge modulaire}
\label{ssec:connexion_jauge}

Pour tout module $m \geq 2$, la récurrence additive
\begin{equation}
  r_{n+1}^{(m)} \;=\; \bigl(r_n^{(m)} + g_n\bigr) \bmod m
  \label{eq:gauge_recursion}
\end{equation}
exprime que la trajectoire des résidus modulaires
$r_n^{(m)} = p_n \bmod m$ est entièrement déterminée par la
donnée initiale et la suite $\{g_n\}$. Cette récurrence est
l'\emph{équation de transport parallèle} associée à la connexion
modulaire de PT : à chaque pas~$n$, le résidu est translaté par
le « gain de jauge » $g_n \bmod m$, et l'on lit la trajectoire
comme un transport le long du « temps arithmétique »~$n$.

Le théorème de connexion de jauge \cite[ch.~1, Thm.~gauge]{PT_Monograph}
formalise cette lecture : la récurrence~\eqref{eq:gauge_recursion}
définit, pour chaque entier $m$ sans facteur carré
$m = q_1 \cdots q_k$ produit de premiers distincts, une connexion
sur le fibré trivial
$\Omega \times (\mathbb{Z}/m\mathbb{Z}) \to \Omega$ au-dessus de
l'espace de chemins $\Omega = \mathbb{N}^{\mathbb{N}}$. Par
décomposition CRT
\begin{equation}
  \mathbb{Z}/m\mathbb{Z}
  \;\cong\;
  \prod_{i=1}^{k} \mathbb{Z}/q_i\mathbb{Z},
  \label{eq:CRT_iso_fr}
\end{equation}
cette connexion se scinde en $k$ composantes indépendantes, une
par premier $q_i$, et l'on peut traiter chaque facteur
$\mathbb{Z}/q_i\mathbb{Z}$ comme un cercle discret de
$q_i$~points portant sa propre connexion de jauge. C'est cette
décomposition qui rend la connexion modulaire de PT structuralement
plus simple qu'une connexion classique sur $\mathbb{R}^4$ : elle
est, par construction, un produit tensoriel d'objets
finis-dimensionnels.

\subsection{L'identité d'holonomie T6}
\label{ssec:T6_holonomie}

Le théorème d'holonomie T6 \cite[ch.~6]{PT_Monograph} donne
l'expression explicite de l'angle d'holonomie associé au transport
parallèle le long d'un cycle de la connexion modulaire. Pour
chaque premier~$p$ et chaque échelle $\mu > 0$,
\begin{align}
  \delta_p(\mu) &\;:=\; \frac{1 - q(\mu)^p}{p},
  &
  \cos \theta_p(\mu) &\;:=\; 1 - \delta_p(\mu),
  \label{eq:delta_theta_fr}
\\
  \sinp{p}(\mu) &\;=\; \delta_p(\mu)\bigl(2 - \delta_p(\mu)\bigr),
  &
  q(\mu) &\;\in\; \{\qplus,\; q_-\},
  \label{eq:sin2_theta_fr}
\end{align}
où $\qplus(\mu) = 1 - 2/\mu$ et $q_-(\mu) = e^{-1/\mu}$ sont les
deux valeurs canoniques du paramètre $q$ singularisées par la
bifurcation de maximum-entropie L0 \cite[ch.~3]{PT_Monograph}. Pour
les questions d'holonomie traitées ici, on adopte la branche
couplage $q = \qplus$, qui décrit le secteur EM et les leptons.

L'identité~\eqref{eq:sin2_theta_fr} est purement algébrique --- elle
ne fait intervenir que la définition trigonométrique
$\sin^2 \theta = 1 - \cos^2 \theta = 1 - (1 - \delta)^2 = \delta
(2 - \delta)$ --- mais elle a une conséquence frappante : la
trigonométrie habituelle, et avec elle la constante $\pi$,
\emph{émerge} de la dynamique modulaire sans avoir été introduite
au départ. Le cercle continu n'est jamais postulé ; il apparaît
comme limite continue du cercle discret $\mathbb{Z}/p\mathbb{Z}$
quand $p \to \infty$, et $\pi$ apparaît via la fonction zêta de
Riemann $\zeta(2) = \pi^2/6 = \prod_p (1 - 1/p^2)^{-1}$, qui
relie directement la mesure du cercle continu au crible des
premiers.

\paragraph{Bilan structural.}
Les trois éléments établis dans cette section dessinent une
architecture dans laquelle la connexion de jauge et l'angle
d'holonomie sont des objets \emph{dérivés}, calculables à partir
du seul champ scalaire $\{g_n\}$, et non des postulats
indépendants. La position C du débat fondationnel ---
l'holonomie comme primitive --- prend, sous PT, un sens
considérablement plus fort qu'elle n'en a dans la théorie des
champs sur variété : elle n'est pas une option d'interprétation,
elle est l'unique sortie d'une construction structurale qui
n'admet aucun autre choix. C'est ce statut que la section
suivante formalise.

\section{L'angle d'holonomie $\theta_p$ comme objet primitif}
\label{sec:holonomie_primitive}

Cette section établit le résultat fondamental qui distingue la
PT des théories classiques de jauge : sous les hypothèses
structurales de la PT, l'angle d'holonomie modulaire $\theta_p$
est l'\emph{unique} invariant primitif satisfaisant
simultanément quatre conditions naturelles, et le couple
classique $(A_\mu, F_{\mu\nu})$ est entièrement reconstructible
comme dérivée d'un objet plus primitif --- le potentiel scalaire
de persistance $S = -\ln \alpha$. La position C du débat
fondationnel n'est plus une interprétation : elle est une
conséquence structurale.

\subsection{Proposition 3.1 : $A_\mu$ et $F_{\mu\nu}$ sont dérivés}
\label{ssec:prop_3_1}

\begin{proposition}[Décomposition du potentiel de persistance]
\label{prop:A_F_derives}
Soit $S(\mu) = -\ln \alpha(\mu)$ le potentiel de persistance. Alors
\begin{itemize}
  \item le 4-potentiel $A_\mu$ est identifiable au gradient
  $\partial_\mu S$ ;
  \item le tenseur de champ $F_{\mu\nu}$ est identifiable aux
  composantes antisymétriques de $\mathrm{Hess}(S)$ ;
  \item la métrique $g_{\mu\nu}$ est identifiable aux composantes
  symétriques de $\mathrm{Hess}(S)$ (Lemme F, ch12 monographie).
\end{itemize}
\end{proposition}

\begin{proof}[Démonstration]
La construction est détaillée au chapitre~12 de la
monographie~\cite{PT_Monograph} ; on en retrace ici les trois
étapes clés.

\textbf{Étape 1 (identification du gradient).}
La métrique de Fisher associée à la distribution stationnaire
$\pi_m$ du crible se calcule comme le Hessien du logarithme de la
fonction de partition $Z_{\rm Ruelle}$, et la fonction
$\alpha(\mu)$ apparaît comme la première fonction propre du
transfert. Par construction, le gradient $\partial_\mu S =
-\partial_\mu \ln \alpha$ joue dans le formalisme PT le rôle que
joue $A_\mu$ dans l'électrodynamique classique : il intervient
linéairement dans le terme d'interaction du Lagrangien effectif
(\S~F3 de la monographie) et il subit, sous les transformations
de jauge modulaire $\theta_p \mapsto \theta_p + 2\pi k$, un
déphasage compatible avec l'invariance de jauge classique.

\textbf{Étape 2 (identification du Hessien antisymétrique).}
Les composantes antisymétriques du Hessien
$\mathrm{Hess}(S)_{\mu\nu} = \partial_\mu \partial_\nu S$ encodent
les courbures dérivées du potentiel scalaire, et reproduisent les
six degrés de liberté de $F_{\mu\nu} = \partial_\mu A_\nu -
\partial_\nu A_\mu$ par identification directe. L'identité
$dA = F$ devient, sous PT, l'identité tautologique
$d(dS) = 0$ projetée sur les composantes antisymétriques de
$\mathrm{Hess}(S)$. Cette projection est non-triviale en présence
de la signature lorentzienne $g_{00} < 0$, qui sépare les
composantes spatiales et temporelles du Hessien.

\textbf{Étape 3 (identification du Hessien symétrique).}
Les composantes symétriques de $\mathrm{Hess}(S)$ reproduisent la
métrique $g_{\mu\nu}$ de l'espace-temps émergent (Bianchi I sur
les premiers actifs $\{3, 5, 7\}$, cf.\ Lemme F du
ch.~12~\cite{PT_Monograph}). Cette identification est, en
elle-même, le contenu du théorème de reconstruction métrique :
la métrique de l'espace-temps de la PT \emph{est} le Hessien du
logarithme de la fonction de partition, sans nouvelle structure
introduite.
\end{proof}

La conséquence immédiate est que $A_\mu, F_{\mu\nu}, g_{\mu\nu}$
sont tous trois \emph{dérivables du même potentiel scalaire}
$S(\mu)$, et que ce potentiel est, en retour, calculé à partir du
seul champ $\{g_n\}$. Aucun de ces trois objets n'est donc
primaire dans l'architecture PT : ils sont tous trois secondaires
par rapport à $S$, et $S$ est lui-même secondaire par rapport à
$\{g_n\}$. C'est sur ce constat structural que repose
l'argument d'unicité de la sous-section suivante.

\subsection{Proposition 3.2 : unicité de $\theta_p$ comme invariant primitif}
\label{ssec:prop_3_2}

\begin{proposition}[Primauté holonomique]
\label{prop:theta_p_unique}
L'angle d'holonomie $\theta_p$ est l'unique invariant primitif
simultanément :
\begin{enumerate}[label=(\roman*),nosep]
  \item jauge-invariant,
  \item sensible à Aharonov--Bohm,
  \item borné,
  \item PT-définissable depuis BA0 seul.
\end{enumerate}
\end{proposition}

\begin{proof}[Démonstration]
Nous montrons par exclusion qu'aucun autre invariant naturel ne
satisfait simultanément les quatre conditions.

\textbf{Exclusion de $A_\mu$.}
Le potentiel $A_\mu$ viole la condition (i) : il dépend de la
jauge par construction. Aucune projection naturelle de $A_\mu$
n'est à la fois jauge-invariante et bornée tout en restant
sensible à Aharonov--Bohm dans la région où $F = 0$. La position
A du débat fondationnel est donc structurellement incompatible
avec les quatre conditions.

\textbf{Exclusion de $F_{\mu\nu}$.}
Le tenseur de champ $F_{\mu\nu}$ satisfait (i), (iii), (iv) mais
viole la condition (ii) : il est, par construction, identiquement
nul dans la région intérieure d'un trajet Aharonov--Bohm idéal,
et ne peut donc pas rendre compte du déphasage observé. La
position B du débat doit ajouter une structure non-locale
supplémentaire pour récupérer la sensibilité à AB ; cette
structure ajoutée n'est pas PT-définissable depuis BA0 seul.

\textbf{Exclusion des fonctions arbitraires de $\{g_n\}$.}
Une fonction arbitraire $F[\{g_n\}]$ peut satisfaire (i), (ii),
(iv) mais viole en général (iii) : sans structure de transport
parallèle modulaire, la fonction n'est pas naturellement bornée
et n'admet pas d'identification simple avec un invariant
topologique mesurable. Seules les fonctions qui se réduisent à
l'angle modulaire $\theta_p$ d'un canal CRT
$\mathbb{Z}/p\mathbb{Z}$ remplissent les quatre conditions
simultanément.

\textbf{Existence et unicité.}
L'angle $\theta_p$ défini par
l'identité~\eqref{eq:sin2_theta_fr} satisfait (i) (il est
jauge-invariant car défini comme intégrale fermée), (ii) (il est
non-nul dans toute région où la connexion modulaire est
non-triviale, y compris dans la région intérieure d'un cycle
AB), (iii) (il prend ses valeurs dans $[0, \pi]$ par
construction), et (iv) (il est défini par
l'identité~\eqref{eq:sin2_theta_fr}, qui ne fait intervenir que
$\{g_n\}$ via $q(\mu)$). C'est l'unique invariant primitif qui
réalise les quatre conditions.
\end{proof}

\subsection{La position Wu--Yang--Healey promue en théorème PT}
\label{ssec:wu_yang_promu}

Les Propositions~\ref{prop:A_F_derives}
et~\ref{prop:theta_p_unique} ont une lecture commune : la
position C du débat fondationnel
(\S~\ref{ssec:position_C}) --- la primauté de l'holonomie ---
n'est plus, dans le cadre PT, un choix d'interprétation parmi
d'autres. Elle est une conséquence structurale : sous les axiomes
de la PT, $\theta_p$ est l'\emph{unique} invariant primitif
disponible, et $A_\mu, F_{\mu\nu}, g_{\mu\nu}$ sont tous trois
dérivés de ce dernier via le potentiel scalaire $S$.

Ce statut a deux conséquences immédiates. D'une part, la PT
\emph{prédit} que toute mesure de $A_\mu$ ou $F_{\mu\nu}$ se
réduit, dans son contenu informationnel, à une mesure de $\theta_p$
sur un canal modulaire actif. Aucune information physique n'est
perdue par cette réduction --- la phase AB est intégralement
encodée dans $\theta_3$ (cf.\ \S~\ref{ssec:ab_consequence}). D'autre
part, la PT \emph{exclut} les ontologies de la position A (la
réalité du potentiel) et de la position B (la non-localité
intrinsèque du champ) : ces ontologies postulent des objets
primitifs qui, dans le cadre PT, ne peuvent exister
indépendamment du potentiel scalaire $S$ d'où elles sont dérivées.

La position C devient ainsi, sous PT, plus qu'une interprétation :
elle est le seul cadre dans lequel les quatre conditions
naturelles (i)--(iv) de la
Proposition~\ref{prop:theta_p_unique} peuvent être satisfaites
simultanément. Le débat fondationnel est, dans cette mesure
précise, tranché.

\section{Le découplage CRT et la non-localité modulaire}
\label{sec:decouplage_crt}

Les Sections~\ref{sec:champ_connexion} et~\ref{sec:holonomie_primitive}
ont établi le statut de l'angle d'holonomie $\theta_p$ comme
invariant primitif et celui des objets classiques $A_\mu,
F_{\mu\nu}, g_{\mu\nu}$ comme dérivés du potentiel scalaire $S$.
Pour relier ce statut à la question expérimentale de la
non-localité, il reste à examiner les corrélations entre canaux
modulaires distincts --- les seules corrélations
``inter-objets'' que la PT autorise à porter une signification
informationnelle non-triviale. Cette section énonce, sans
démonstration détaillée, les deux théorèmes structurels prouvés
dans l'article companion technique~\cite{PT_CRT_DECOUPLING}
ainsi que leur corollaire empirique.

\subsection{Théorème 4.1 : Théorème de découplage géométrique (forme empirique)}
\label{ssec:thm_4_1_fr}

\begin{theoreme}[Théorème de découplage géométrique, forme empirique]
\label{thm:invariance_geom}
\footnote{Vérifié formellement en Lean~4 + Mathlib comme
\texttt{PT.CrtDecoupling.Empirical.empirical\_invariance}
(forme hypothétique),
\texttt{SpectralReduction.empirical\_invariance\_closed}
(forme fermée fini),
\texttt{Phase3.empirical\_invariance\_infinite}
(forme infinie via Birkhoff--Hopf), et
\texttt{Phase4.sieve\_empirical\_invariance}
(clôture spécifique au crible) ;
cf.~\cite[\S~8.5]{PT_CRT_DECOUPLING}.}
Soient $p, q \in \{3, 5, 7\}$ premiers actifs distincts, et soient
$\Phi_p \in \mathcal{A}_p$, $\Phi_q \in \mathcal{A}_q$ des
observables bornées sur les $\sigma$-algèbres modulaires
correspondantes. La fonction
\begin{equation}
  \mathcal{I}_{p,q}(\mu)
  \;:=\;
  I_{\pi_m^{\mathrm{emp}}}\!\bigl(\Phi_p\,;\,\Phi_q \,\big|\, G(\mu)\bigr)
\end{equation}
est constante sur $\mu \in [\mu_c, \infty)$.
\end{theoreme}

La démonstration complète est donnée dans le théorème~6.1
du companion technique \cite{PT_CRT_DECOUPLING}. Elle procède en
quatre étapes : (i) la mesure empirique $\pi_m^{\rm emp}$ de la
trajectoire des résidus est une fonctionnelle du seul champ
$\{g_n\}$, donc indépendante de $\mu$ ; (ii) la géométrie
$G(\mu)$ est mesurable par rapport à la tribu engendrée par le
spectre de la matrice de transfert $T_m$, qui est elle-même
contenue dans la tribu invariante du shift ; (iii) la dynamique du
crible est ergodique par le théorème de Birkhoff--Hopf appliqué
au gap spectral T4, donc la tribu invariante est triviale
modulo $\pi_m^{\rm emp}$, donc $G(\mu)$ est presque sûrement
constante ; (iv) conditionner par une variable presque sûrement
constante ne change pas l'information mutuelle. La conclusion
suit immédiatement.

Le point conceptuel est que la métrique émergente $G(\mu)$ --- y
compris à travers la transition de Hartle--Hawking en $\mu_c
\approx 6{,}97$ et jusqu'au point fixe canonique
$\mustar = 15$ --- ne modifie jamais le contenu informationnel
joint des canaux modulaires actifs. Les corrélations
inter-canaux sont \emph{intrinsèquement} pré-géométriques : elles
existent avant la construction de la métrique et lui sont
indifférentes.

\subsection{Théorème 4.2 : Théorème de découplage géométrique (forme Ruelle)}
\label{ssec:thm_4_2_fr}

\begin{theoreme}[Théorème de découplage géométrique, forme Ruelle]
\label{thm:ruelle_zero_fr}
\footnote{Vérifié formellement en Lean~4 + Mathlib comme
\texttt{PT.CrtDecoupling.Main.geometric\_decoupling\_ruelle} ;
cf.~\cite[\S~8.5]{PT_CRT_DECOUPLING}.}
Sous la mesure idéalisée $\pi_m^{\mathrm{Ruelle}}$, on a
$I(\Phi_p\,;\,\Phi_q) = 0$ exactement.
\end{theoreme}

La démonstration est immédiate à partir de la factorisation
tensorielle CRT du transfert
$T_{pq} = T_p \otimes T_q$ (Théorème~3.1
de~\cite{PT_CRT_DECOUPLING}) et de la factorisation correspondante
de l'unique vecteur propre gauche de Perron
$\pi_{pq}^{\rm Ruelle} = \pi_p^{\rm Ruelle} \otimes \pi_q^{\rm Ruelle}$
(Lemme~3.2 du même article). Cette factorisation est une
condition d'indépendance, donc la divergence de Kullback--Leibler
entre la loi jointe et le produit des marginales est nulle ;
par l'inégalité de traitement de données, l'information mutuelle
de toute paire d'observables bornées en hérite.

\emph{Validation numérique.} La factorisation
$\pi_{15}^{\rm Ruelle} = \pi_3^{\rm Ruelle} \otimes \pi_5^{\rm Ruelle}$
est vérifiée au sens du résidu
\begin{equation}
  \|\pi_{15}^{\rm direct} - \pi_3 \otimes \pi_5\|_2
  \;=\; 5{,}88 \times 10^{-53}
\end{equation}
à \texttt{mpmath} 50~chiffres de précision (test T2 PASS de
\texttt{scripts/test\_crt\_decoupling.py}~\cite{PT_CRT_DECOUPLING}).

\subsection{Corollaire 4.3 : Forme fermée du découplage}
\label{ssec:cor_4_3_fr}

\begin{corollaire}[Forme fermée du découplage]
\label{cor:closed_form_fr}
\footnote{Non formalisable comme théorème Lean en l'état car le
paramètre de profondeur $k_{p,q}^{\rm eff}$ est ajusté
empiriquement ; seule la prédiction à $k$ fixé admet un énoncé
Lean-vérifiable ;
cf.~\cite[\S~8.5]{PT_CRT_DECOUPLING}.}
La constante $\mathcal{I}_{p,q}$ admet une expression fermée
dérivée des angles d'holonomie :
\begin{equation}
  \mathcal{I}_{p,q}
  \;=\;
  \frac{1}{2 \ln 2}\Bigl[
    \cos^{2 k_{p,q}}(\theta_3, \mustar)
    + \sin^{4 k_{p,q}}(\theta_2, p_1)
  \Bigr].
  \label{eq:closed_form_fr}
\end{equation}
Validations numériques (cf.~\cite{PT_CRT_DECOUPLING}, Tableau B) :
$\mathcal{I}_{3,5} = 0{,}132$ bits (cible $0{,}138$, écart $4{,}5\%$) ;
$\mathcal{I}_{3,7} = 0{,}306$ bits (cible $0{,}283$, écart $8{,}1\%$).
\end{corollaire}

Le paramètre de profondeur $k_{p,q}$ est ajusté empiriquement à
$k_{3,5}^{\rm eff} = 7$ et $k_{3,7}^{\rm eff} = 4$ ; une dérivation
combinatoire à partir de la structure antidiagonale de $T_3$ reste
ouverte (cf.\ remarque~7.1 de~\cite{PT_CRT_DECOUPLING}). Cette
incertitude ne porte que sur la \emph{valeur quantitative} de
$\mathcal{I}_{p,q}$ ; le résultat structurel
(Théorème~\ref{thm:invariance_geom}) ne dépend pas du choix de
$k_{p,q}$.

\subsection{La non-localité spatiale dissoute en non-localité modulaire}
\label{ssec:dissolution}

Les trois énoncés précédents permettent de reformuler en termes
PT le statut de la « non-localité » apparente des corrélations
quantiques. On peut résumer la situation en trois propositions
strictement parallèles aux trois positions du débat fondationnel
(\S~\ref{sec:debat}), mais avec un statut différent :

\paragraph{(i) Les corrélations inter-canaux existent.}
$\mathcal{I}_{3,5}$ et $\mathcal{I}_{3,7}$ sont
non nuls et de magnitude prédite par le
Corollaire~\ref{cor:closed_form_fr}. La PT ne nie pas l'existence
des corrélations EPR/Bell ; elle en fournit, au contraire, une
prédiction quantitative.

\paragraph{(ii) Ces corrélations sont d'origine arithmétique.}
La forme fermée du Corollaire~\ref{cor:closed_form_fr}
fait intervenir uniquement les angles d'holonomie $\theta_3,
\theta_2$ évalués aux échelles canoniques $\mustar = 15$ et
$p_1 = 3$. Aucune intégrale spatiale, aucune métrique, aucune
distance entre points de l'espace n'entre dans l'expression. La
non-localité observée n'est pas une non-localité dans
l'espace : c'est une corrélation intra-CRT entre canaux modulaires
différents.

\paragraph{(iii) Ces corrélations sont invariantes par $G(\mu)$.}
Le Théorème~\ref{thm:invariance_geom} interdit que la géométrie
émergente \emph{couple} les canaux modulaires. Quelle que soit la
position dans la famille géométrique
$\mu \in [\mu_c, \infty)$ --- y compris à toute distance entre
détecteurs dans un dispositif Bell --- l'information mutuelle entre
canaux distincts reste la même.

Ces trois propositions ensemble dissolvent la non-localité
spatiale au sens fort de la position B : les corrélations
quantiques ne traversent pas l'espace parce qu'elles ne sont pas
dans l'espace. Elles sont dans la structure CRT, qui se manifeste
dans la géométrie émergente $G$ par projection, sans la traverser
ni l'utiliser comme support causal. Le paradoxe apparent d'une
information « plus rapide que $c$ » dans les tests de Bell est un
artéfact de la lecture spatiale d'une corrélation qui, dans la
substrat PT, n'est pas spatiale. Nous revenons sur les conséquences
expérimentales de cette dissolution dans
la~\S~\ref{ssec:bell_consequence}.

\section{Prédiction QH1a falsifiable}
\label{sec:prediction_qh1a}

Les Sections~\ref{sec:champ_connexion}--\ref{sec:decouplage_crt}
établissent la position structurelle de la PT par rapport au
débat fondationnel : l'angle d'holonomie est primitif, les
objets classiques sont dérivés, les corrélations entre canaux
modulaires sont arithmétiques et invariantes par la géométrie
émergente. Cette position n'est pas qu'interprétative : elle
implique une prédiction expérimentale précise, falsifiable sur
plusieurs plateformes existantes, et structuralement
distinguable de toute prédiction d'une physique standard sur
variété courbe.

\subsection{Énoncé de la prédiction}
\label{ssec:qh1a_enonce}

\begin{prediction}[QH1a : invariance géométrique de la MI]
\label{pred:qh1a}
Dans un système atomique multi-électronique sondé par
spectroscopie CRT (canaux $p = 3, 5, 7$), l'information mutuelle
$I(n^{(p)}\,;\,n^{(q)})$ entre nombres d'occupation modulaires :
\begin{enumerate}[label=(\roman*),nosep]
  \item est non-nulle et donnée par la formule fermée du
  Corollaire~\ref{cor:closed_form_fr} (test arithmétique) ;
  \item est invariante sous toute transformation géométrique
  du système (déplacement nucléaire, déformation de la maille
  cristalline, séparation des détecteurs) (test géométrique).
\end{enumerate}
\end{prediction}

La condition (i) est le contenu quantitatif de
Corollaire~\ref{cor:closed_form_fr} : la MI prend les valeurs
$\mathcal{I}_{3,5} = 0{,}138$ bits et
$\mathcal{I}_{3,7} = 0{,}283$ bits (valeurs corpus), prédictibles
au $\pm 20\%$ près par la formule fermée à un seul paramètre
entier $k_{p,q}$ par paire. Toute mesure qui ne reproduit pas ces
valeurs (à la précision statistique de l'expérience) invalide la
prédiction.

La condition (ii) est le contenu structurel de
Théorème~\ref{thm:invariance_geom} : la MI doit rester constante
sous toute reparamétrisation de la géométrie spatiale. Toute
mesure qui détecte une dépendance de la MI en la séparation des
détecteurs, la distance internucléaire, ou tout autre paramètre
géométrique, invalide la prédiction.

\subsection{Critère de falsification}
\label{ssec:qh1a_falsification}

Soit $G_0, G_1$ deux configurations géométriques distinctes d'un
même système quantique, et soient
$\mathcal{I}_{p,q}(G_0), \mathcal{I}_{p,q}(G_1)$ les valeurs de
l'information mutuelle mesurée sous chacune. Définissons
\begin{equation}
  \Delta\mathcal{I}_{p,q}
  \;:=\;
  \mathcal{I}_{p,q}(G_1) - \mathcal{I}_{p,q}(G_0)
  \label{eq:Delta_I_def}
\end{equation}
et soit $\epsilon_{\rm stat}$ l'incertitude statistique de
$\Delta\mathcal{I}_{p,q}$ telle qu'évaluée par la procédure de
mesure. La PT prédit
\begin{equation}
  \bigl|\Delta\mathcal{I}_{p,q}\bigr| \;<\; \epsilon_{\rm stat}
  \qquad \text{pour toute paire } (G_0, G_1).
  \label{eq:QH1a_predict}
\end{equation}
\emph{Critère de falsification.} Si une expérience mesure
$|\Delta\mathcal{I}_{p,q}| > 5 \epsilon_{\rm stat}$ pour une
paire de configurations géométriques distinctes, la prédiction
QH1a est falsifiée, et avec elle le théorème
d'invariance~\ref{thm:invariance_geom}. Une telle falsification
serait également une réfutation directe de la PT, puisque
$\theta_p$ et $\mathcal{A}_p$ sont des objets structuraux de la
théorie et ne peuvent être ajustés pour absorber un signal
géométrique.

\subsection{Plateformes expérimentales candidates}
\label{ssec:qh1a_plateformes}

Trois plateformes contemporaines offrent l'accès aux nombres
d'occupation modulaires $n^{(p)}$ avec contrôle indépendant de la
géométrie spatiale.

\paragraph{(a) Atomes froids dans réseaux optiques.}
Un réseau optique à espacement variable permet de scanner la
distance entre puits de potentiel adjacents tout en mesurant la
distribution d'occupation modulaire des niveaux internes par
spectroscopie de fluorescence. La géométrie est contrôlée par la
longueur d'onde du laser de piégeage ; les nombres d'occupation
$n^{(p)}$ sont accessibles par mesure résonante au niveau du
résidu modulaire. Précision attendue : $\epsilon_{\rm stat}
\sim 10^{-3}$ bits sur $\mathcal{I}_{3,5}$ pour
$N \sim 10^4$ atomes piégés, soit une sensibilité largement
suffisante pour falsifier QH1a à $5\sigma$ en présence d'un
couplage géométrique.

\paragraph{(b) Ions piégés à séparations contrôlées.}
Une chaîne d'ions Yb$^+$ ou Be$^+$ permet de varier la séparation
inter-ion avec une précision micrométrique, tout en mesurant les
corrélations inter-canaux par tomographie d'état. La structure
hyperfine des niveaux internes fournit les nombres d'occupation
modulaires pour $p = 3$ (canal de couleur SU(3) effectif) et
$p = 5$ (canal de saveur dans les configurations à plusieurs
électrons de valence). Précision attendue : $\epsilon_{\rm stat}
\sim 10^{-4}$ bits, soit la plateforme la plus contraignante.

\paragraph{(c) Molécules diatomiques à distance internucléaire ajustable.}
Une molécule diatomique sondée par spectroscopie laser à
résolution sub-MHz permet de varier la distance internucléaire
$R$ via un champ électrique externe (effet Stark) et de mesurer
simultanément les nombres d'occupation modulaires des niveaux
électroniques. La signature recherchée est l'absence de
dépendance de $\mathcal{I}_{3,5}$ en $R$ sur une plage de
plusieurs angstroms. Précision attendue : $\epsilon_{\rm stat}
\sim 5 \times 10^{-3}$ bits, plus modeste mais structuralement
complémentaire (le couplage géométrique attendu d'une théorie
standard est ici dominé par la variation du potentiel
vibrationnel).

\subsection{Contraste avec la physique standard}
\label{ssec:qh1a_contraste}

Une théorie quantique standard sur variété courbe prédit
génériquement $|d\mathcal{I}_{p,q}/dG| \neq 0$. La raison
structurelle est que les champs de matière sont, dans le
formalisme standard, des sections de fibrés au-dessus de la
variété courbe, et leurs corrélations héritent de la structure
métrique de la base par les couplages gravitationnels minimaux.
L'effet est typiquement faible (suppression par $1/M_{\rm Pl}^2$
ou par le couplage gravitationnel local), mais il est
\emph{non-nul} et calculable.

La PT, au contraire, prédit $d\mathcal{I}_{p,q}/dG = 0$
identiquement. La raison est que les canaux modulaires
$\mathcal{A}_p$ sont définis avant la métrique $G$ et lui sont
mesurablement antérieurs. Aucune fonction de la métrique
n'apparaît dans la formule
fermée~\eqref{eq:closed_form_fr} ; la dérivée par rapport à $G$
est exactement zéro.

Ce contraste rend QH1a un test discriminant fort entre les deux
cadres. Une mesure positive ($|d\mathcal{I}/dG| > 5\epsilon$)
falsifie la PT et confirme le cadre standard ; une mesure nulle
($|d\mathcal{I}/dG| < \epsilon$) confirme la PT au niveau
structurel et exclut tout couplage géométrique minimal au-dessus
de la précision atteinte. Aucune des plateformes actuelles n'a
encore réalisé la mesure complète --- en particulier, aucune
n'a scanné la dépendance géométrique de la MI sur plusieurs ordres
de grandeur de séparation. C'est la principale lacune
expérimentale à combler pour soumettre la PT à un test
falsificationniste réel.

\section{Conséquences pour Aharonov--Bohm, Bell, EPR}
\label{sec:consequences_bell_ab_epr}

Les trois résultats structurels précédents
--- l'angle d'holonomie comme primitive
(Proposition~\ref{prop:theta_p_unique}), le découplage géométrique
(Théorème~\ref{thm:invariance_geom}), et la forme fermée
arithmétique des corrélations
(Corollaire~\ref{cor:closed_form_fr}) --- modifient la lecture
de trois expériences classiques de la physique fondamentale : la
configuration Aharonov--Bohm, les tests de Bell, et la
configuration EPR originelle. Nous examinons brièvement chacune.

\subsection{Aharonov--Bohm : la phase est $\theta_3$, pas $A_\mu$}
\label{ssec:ab_consequence}

Dans la lecture PT, le déphasage interférométrique
$\Delta\varphi$ observé dans une configuration
Aharonov--Bohm \cite{AharonovBohm1959} n'est pas
$(e/\hbar) \oint A \cdot d\ell$ mais $\theta_3$, l'angle
d'holonomie du canal modulaire $p = 3$ associé au flux entouré
par le trajet de la particule. L'identification entre les deux
formules est tautologique en présence du dictionnaire DIC qui
traduit les unités PT en unités SI ; ce qui distingue la lecture
PT n'est pas la valeur numérique du déphasage mais l'identification
du \emph{support physique} du déphasage.

Cette identification a deux conséquences. D'abord, la phase AB
est un invariant strictement topologique du chemin~$\gamma$ :
elle ne dépend que de la classe d'homotopie de $\gamma$ dans le
complémentaire de la région de flux, et est insensible à toute
déformation continue de $\gamma$ dans cette classe. Cette
topologie n'est pas une propriété de la métrique $G(\mu)$ ; elle
est antérieure à la métrique, dans le sens du
Théorème~\ref{thm:invariance_geom}. Ensuite, l'absence de champ
$F = 0$ dans la région intérieure d'un solénoïde idéal n'est plus
un paradoxe : $F_{\mu\nu}$ n'est pas l'objet qui porte
l'information de la phase. C'est $\theta_3$ qui la porte, et
$\theta_3$ est non-nul même là où $F$ s'annule, parce que
$\theta_3$ est défini par la connexion modulaire de PT, qui
n'est elle-même pas une fonction locale du champ.

\subsection{EPR/Bell : corrélations intra-canal, pas inter-canal}
\label{ssec:bell_consequence}

Les tests de Bell \cite{Bell1964,Aspect1982,Hensen2015}
établissent que les corrélations observées entre détecteurs
spatialement séparés excèdent la borne de la
CHSH \cite{CHSH1969} attendue d'une théorie locale à variables
cachées. Dans la lecture standard, ces corrélations sont
interprétées comme la signature d'une non-localité fondamentale du
champ quantique : l'effondrement d'un état EPR sur une mesure à
gauche détermine instantanément l'état corrélé à droite, à toute
distance, et apparemment plus vite que la lumière (mais sans
permettre de communication superluminale, conformément au
no-signalling).

Dans la lecture PT, cette interprétation se renverse. Les deux
mesures effectuées dans un test de Bell ne portent \emph{pas} sur
deux canaux modulaires distincts (auquel cas, par
Théorème~\ref{thm:invariance_geom}, leur MI serait invariante par
géométrie et de magnitude prédite par la formule fermée
arithmétique~\eqref{cor:closed_form_fr}). Elles portent sur le
\emph{même} canal modulaire --- typiquement $p = 2$ pour les
expériences de polarisation/spin, ou $p = 3$ pour les expériences
de couleur de quarks (analogues). Une corrélation intra-canal
n'est pas l'objet du
Théorème~\ref{thm:invariance_geom} : c'est une corrélation
d'ordre supérieur dans la structure CRT, indépendante de toute
notion de séparation spatiale, parce que le canal modulaire
existe avant la construction de l'espace par projection.

Cette lecture s'accorde avec la signature
empirique des tests les plus précis : la corrélation observée
entre deux détecteurs séparés de plusieurs kilomètres
\cite{Hensen2015} est strictement identique à celle observée à
quelques mètres \cite{Aspect1982}, à la précision des mesures
près. Aucune dépendance significative en la distance n'a jamais
été détectée. Cette absence de dépendance est, dans la lecture
PT, l'attendu structurel : ce qui est mesuré n'est pas dans
l'espace, donc la séparation spatiale ne peut pas modifier le
résultat.

\subsection{Dissolution de la non-localité spatiale}
\label{ssec:dissolution_consequence}

La conséquence ontologique des deux sous-sections précédentes est
la \emph{dissolution} (statut [DISS]) de la non-localité spatiale
au sens fort. Le « paradoxe » d'une information apparemment
transmise plus vite que $c$ disparaît parce que les corrélations
ne se propagent pas dans l'espace : elles existent dans la
structure CRT, et l'espace n'est qu'une projection émergente de
cette structure via la métrique $G(\mu)$.

Le statut [DISS] est, dans la nomenclature PT, plus faible que
[THM] mais plus fort que la simple ré-interprétation. Une
dissolution affirme que le problème, tel qu'il était posé dans le
cadre classique, est mal-posé : ses prémisses
implicites --- l'existence d'un espace classique de fond et la
nécessité que toute corrélation observée s'y propage --- ne sont
pas valides en PT. Le problème ne se résout donc pas dans le sens
où l'on aurait trouvé la « vraie » explication classique des
corrélations EPR ; il se dissout dans le sens où les corrélations
sont vues comme des objets qui n'appartenaient jamais au domaine
spatial.

L'horizon expérimental est alors clair : la PT ne prédit aucun
signal qui puisse séparer deux détecteurs spatialement séparés
au-delà de la corrélation arithmétique CRT, et exclut
catégoriquement tout signal qui dépendrait de la métrique
$G(\mu)$ entre les deux détecteurs. Toute observation positive
d'une telle dépendance falsifie simultanément la PT et la
lecture présentée ici. Réciproquement, la confirmation
expérimentale de l'invariance géométrique des corrélations EPR ---
i.e.\ la prédiction QH1a appliquée au régime relativiste --- serait
une validation directe de la lecture PT du débat fondationnel.

\section{Trois lectures ontologiques (R1, R2, R3)}
\label{sec:lectures_ontologiques}

La PT, par sa nature de théorie purement informationnelle dérivée
d'un unique paramètre arithmétique $s = 1/2$, admet trois lectures
ontologiques compatibles avec sa structure formelle. Cette
trichotomie est exposée en détail dans les chapitres~24 et~26 de
la monographie~\cite{PT_Monograph} ; nous la reprenons ici
brièvement, et l'appliquons spécifiquement au débat
Aharonov--Bohm / Bell développé dans cet article. Le point
crucial à souligner --- et qui constitue l'essentiel
de~\S~\ref{ssec:independance} --- est que les démonstrations des
sections~\ref{sec:champ_connexion}--\ref{sec:prediction_qh1a}
sont logiquement indépendantes du choix de lecture. La trichotomie
porte sur l'\emph{interprétation} de ce que les théorèmes
disent, non sur leur validité.

\subsection{R1 : la PT comme coïncidence mathématique}
\label{ssec:R1}

La lecture R1 considère que la PT reproduit le bon contenu
empirique --- la phase Aharonov--Bohm, les corrélations EPR,
l'absence de dépendance géométrique de la MI --- par une
coïncidence structurale entre la combinatoire du crible
d'Eratosthène et la structure phénoménologique de
l'électrodynamique quantique. Aucun engagement ontologique n'est
pris sur la nature physique des angles $\theta_p$ : ils sont des
quantités mathématiques utiles à la prédiction, sans plus. Les
théorèmes des sections précédentes restent valides --- ils sont
des assertions mathématiques sur la structure du crible --- mais
leur correspondance avec les phénomènes observés est traitée comme
un fait empirique inexpliqué.

Cette lecture est la plus prudente épistémologiquement : elle
n'engage que le minimum nécessaire pour rendre compte des
observations. Son point faible est qu'elle ne fournit aucune
\emph{explication} de la coïncidence : pourquoi la combinatoire
du crible donne-t-elle exactement la bonne structure
phénoménologique ? La lecture R1 ne répond pas à cette question.

\subsection{R2 : structuralisme}
\label{ssec:R2}

La lecture R2 considère que le crible d'Eratosthène et
l'électrodynamique quantique partagent un \emph{pattern}
holonomique abstrait, et que ni l'un ni l'autre n'est
ontologiquement premier. Les $\theta_p$ ne sont alors
ni des objets mathématiques (au sens R1) ni des objets physiques
(au sens R3) : ils sont les composantes d'une structure abstraite
dont les deux théories sont des incarnations.

Cette lecture, classique dans la philosophie des sciences sous
le nom de \emph{structuralisme} \cite{Belot1998,Maudlin2011},
trouve en PT un terrain particulièrement favorable : la
correspondance entre le crible et la phénoménologie n'est plus
une coïncidence mais l'expression d'une structure commune
sous-jacente, et la primauté de l'holonomie devient une propriété
structurelle plutôt qu'une question d'ontologie matérielle. Les
théorèmes des sections précédentes conservent leur validité, et
acquièrent un statut explicatif au sens où ils décrivent la
structure commune.

\subsection{R3 : réalisme informationnel}
\label{ssec:R3}

La lecture R3 considère que les $\theta_p$ \emph{sont} la réalité
physique fondamentale, et que les objets classiques de
l'électrodynamique --- $A_\mu$, $F_{\mu\nu}$, $g_{\mu\nu}$ ---
sont des projections phénoménologiques sur l'espace-temps
émergent. Sous cette lecture, la substance ontologique de la
théorie est l'information ($\DKL$ en bits) ; les notions de
particule, champ, espace, et temps sont des constructions
émergentes au-dessus de cette substance informationnelle.

R3 est la lecture la plus engagée : elle affirme que le monde est
fait d'information arithmétique, que la matière et l'espace n'en
sont que des projections, et que les corrélations observées dans
les tests de Bell sont des reflets directs de la structure CRT
fondamentale. Sous R3, la dissolution de la non-localité spatiale
(\S~\ref{ssec:dissolution_consequence}) prend son sens le plus
fort : les corrélations EPR sont littéralement non-spatiales,
parce que ce qui est mesuré n'a pas de support spatial.

R3 hérite des difficultés philosophiques classiques du réalisme
informationnel \cite{Healey2007} : la nature de
l'« exécution » du crible reste ouverte (qui ou quoi
« calcule » la suite des $\{g_n\}$?), et l'apparente
correspondance entre la structure mathématique et la
phénoménologie observée demande, sous R3, une explication
métaphysique qui dépasse la portée de la théorie elle-même.

\subsection{Indépendance des théorèmes par rapport à la lecture}
\label{ssec:independance}

Le point logique central de cette section est que les
Théorèmes~\ref{thm:invariance_geom}
et~\ref{thm:ruelle_zero_fr}, le
Corollaire~\ref{cor:closed_form_fr}, et la
Prédiction~\ref{pred:qh1a} sont \emph{tous} valides sous R1, R2,
R3. Aucun résultat formel des sections
précédentes ne dépend du choix d'interprétation ontologique. Les
démonstrations sont mathématiques ; les prédictions sont des
fonctions des angles d'holonomie et de la métrique émergente ; et
les valeurs numériques (par exemple
$\mathcal{I}_{3,5} = 0{,}138$ bits) sont calculables
indépendamment de la lecture choisie.

Ce qui change avec la lecture, c'est uniquement
l'\emph{interprétation} de ce que mesure l'expérience. Sous R1,
une confirmation expérimentale de QH1a est un fait empirique
inexpliqué de coïncidence avec la structure du crible. Sous R2,
elle est l'expression empirique d'une structure abstraite
commune entre crible et électrodynamique. Sous R3, elle est une
mesure directe de la réalité ontologique informationnelle. Les
trois lectures sont compatibles avec les mêmes observations ; le
choix entre elles dépend de critères extra-empiriques (parcimonie,
pouvoir explicatif, engagement métaphysique).

Cette indépendance a une conséquence pratique importante pour
l'évaluation expérimentale de la PT : un test de QH1a peut être
mené, et ses résultats interprétés, sans engagement préalable sur
le statut ontologique des $\theta_p$. La théorie est testable
indépendamment de la lecture. C'est, à notre connaissance, le
premier cadre théorique qui propose une dissolution structurelle
du débat Aharonov--Bohm / Bell tout en restant
empiriquement falsifiable, et qui peut être évalué sans avoir à
choisir d'abord entre les positions ontologiques classiques.


\bibliographystyle{plain}
\bibliography{references}

\end{document}
